% team-meeting-2026-10-27-slides.tex -- discussion slides for the team meeting
% on the finished scenario runs, 2026-10-27.
% Build: plain pdflatex (emacs, emacs-r container) or xelatex
% (bittensor-ubuntu), run twice. beamerthemeozlabs.sty, ozcolors.sty and
% ozfonts.sty in this directory are symlinks to ~/code/oz-templates/, and the
% logo is set by absolute path below, so no TEXINPUTS setting is needed.
% Numbers come from output/review-workflow-b-2026-10-26/ (paired_headline.txt,
% model_sets.txt, wb_check.txt) and docs/reviews/workflow-b-full-run-review-2026-10-26.org.
\documentclass[aspectratio=169,11pt]{beamer}

\usetheme{ozlabs}

\renewcommand{\ozfootbrand}{OZ Labs}
\renewcommand{\ozfooturl}{ozstats.space}
\renewcommand{\ozeyelogo}{/var/home/aoz/code/oz-templates/assets/ozlabs-seed-logo.pdf}

\usepackage{booktabs}

\title{Scenario Runs Complete: Choices and First Results}
\subtitle{Heat and cold deaths, 204 countries, 2022--2050, four climate scenarios}
\author{Aaron Osgood-Zimmerman}
\institute{OZ Labs LLC, for the World Bank climate-health project}
\date{October 27, 2026}

\begin{document}

\begin{frame}[plain,noframenumbering]
  \titlepage
\end{frame}

\begin{frame}{Where we are}
  \begin{itemize}
    \item All four scenarios (SSP1-2.6, SSP2-4.5, SSP3-7.0, SSP5-8.5) have
      finished for all 204 countries, 2022--2050, along with the step that
      adjusts each scenario to IHME's mortality forecasts.
    \item I have checked everything Caspar sent. It all passes: the totals
      add up, nothing is missing, and the results match the earlier test
      runs exactly.
    \item Today: four choices I'd like the team's view on, then four early
      results.
    \item \textbf{Caveat for every number here:} these are averages across
      climate models, global totals, with no uncertainty ranges yet. The
      ranges come with the combined summary (choice 4).
  \end{itemize}
\end{frame}

\begin{frame}{Choice 1: which climate models go into the headline numbers?}
  \small
  The adjustment step needs each model's SSP2-4.5 run, and two models don't
  have one. So we have 26 models for SSP1-2.6, 24 for SSP3-7.0 and 27 for
  SSP5-8.5. Only 24 models were run for all four scenarios.

  \medskip
  Heat deaths in 2041--2050, relative to SSP2-4.5:

  \medskip
  \centering
  \begin{tabular}{lccc}
    \toprule
    & Every model & Same models in both & Same 24 everywhere \\
    \midrule
    SSP1-2.6 & 0.985 & 0.963 & 0.964 \\
    SSP3-7.0 & 1.002 & 1.008 & 1.008 \\
    SSP5-8.5 & 1.152 & 1.134 & 1.127 \\
    \bottomrule
  \end{tabular}

  \medskip
  \raggedright
  \begin{itemize}
    \item ``Every model'' compares different sets of models, which adds about
      2 points to the SSP5-8.5 increase. Two of the extra models are among the
      three hottest.
    \item \textbf{My plan:} same models in both for each scenario against
      SSP2-4.5, and the common 24 for any table showing all four side by side.
  \end{itemize}
\end{frame}

\begin{frame}{Choice 2: keep the adjustment to IHME's forecasts?}
  \begin{itemize}
    \item The adjustment barely moves the results: global deaths change by
      less than 0.15\% in every scenario and decade.
    \item For a single country and year it stays within 5\% in 99.9\% of
      cases. The largest effect is Singapore, up to 16.5\%.
    \item The cost is losing one or two climate models per scenario
      (choice 1).
    \item \textbf{My suggestion:} keep the adjusted numbers as the headline,
      since that is the method we agreed, and note that the unadjusted
      numbers agree.
  \end{itemize}
\end{frame}

\begin{frame}{Choice 3: how to present heat and cold}
  SSP5-8.5 relative to SSP2-4.5, 2041--2050:

  \medskip
  \centering
  \begin{tabular}{lcc}
    \toprule
    & Change & Range across models \\
    \midrule
    Heat deaths & +13.4\% & +4\% to +25\% \\
    Cold deaths & $-2.6$\% & $-5.5$\% to $-1.1$\% \\
    All temperature deaths & +1.8\% & $-0.3$\% to +4\% \\
    \bottomrule
  \end{tabular}

  \medskip
  \raggedright
  \begin{itemize}
    \item Cold deaths are about 2.7 times heat deaths, so the drop in cold
      deaths largely offsets the rise in heat deaths.
    \item Reporting only the total would hide the heat increase.
    \item \textbf{My suggestion:} always show heat, cold and the total side by
      side.
  \end{itemize}
\end{frame}

\begin{frame}{Choice 4: what should the combined summary contain?}
  \begin{itemize}
    \item The combined summary pools all models and gives each number an
      uncertainty range that includes disagreement between climate models
      (agreed in September).
    \item It needs one more run on Caspar's machine. The code is written and
      tested. I'll ask him for every year, so we can show any period later
      (single years, five or ten years) without asking him again.
    \item Questions for the team (none of these hold up Caspar's run):
      \begin{enumerate}
        \item \textbf{Periods:} which years or periods should the tables show?
        \item \textbf{Regions:} GBD regions, World Bank regions, or income
          groups?
        \item \textbf{Measures:} deaths and years of life lost; are rates per
          100,000 people also wanted?
      \end{enumerate}
  \end{itemize}
\end{frame}

\begin{frame}{Result 1: most of the growth to 2050 happens in every scenario}
  \small
  Global heat deaths per year, in millions:

  \medskip
  \centering
  \begin{tabular}{lccc}
    \toprule
    & 2022--2030 & 2031--2040 & 2041--2050 \\
    \midrule
    SSP1-2.6 & 0.44 & 0.57 & 0.71 \\
    SSP2-4.5 & 0.44 & 0.56 & 0.74 \\
    SSP3-7.0 & 0.42 & 0.56 & 0.74 \\
    SSP5-8.5 & 0.44 & 0.59 & 0.84 \\
    \bottomrule
  \end{tabular}

  \medskip
  \raggedright
  \begin{itemize}
    \item Heat deaths rise by about 70\% by 2041--2050 even under SSP2-4.5.
      The gap between scenarios is much smaller.
    \item My guess is population aging, IHME's mortality trends and warming
      that is already locked in. I haven't separated these.
    \item Question: do we lead with the overall rise or with the differences
      between scenarios?
  \end{itemize}
\end{frame}

\begin{frame}{Result 2: SSP1-2.6 is slightly worse than SSP2-4.5 at first}
  \begin{itemize}
    \item By 2041--2050 the scenarios are in the expected order: SSP5-8.5
      highest, SSP3-7.0 about the same as SSP2-4.5, SSP1-2.6 lowest.
    \item Before that, SSP1-2.6 has slightly more heat deaths than SSP2-4.5:
      2\% more in 2022--2030 and 1.6\% more in 2031--2040.
    \item The temperatures agree. Where most deaths occur, SSP1-2.6 is
      0.11\,\textcelsius{} warmer than SSP2-4.5 in 2030 and
      0.25\,\textcelsius{} cooler by 2050.
    \item One possible reason is faster cuts in air pollution under SSP1-2.6,
      which removes some cooling haze. I haven't verified this.
    \item Question: does anyone know this pattern from the climate
      literature? Explain it, or just footnote it?
  \end{itemize}
\end{frame}

\begin{frame}{Result 3: SSP3-7.0 gives fewer heat deaths in parts of Europe}
  \small
  SSP3-7.0 relative to SSP2-4.5 (24 models):

  \medskip
  \centering
  \begin{tabular}{lccc}
    \toprule
    & Heat deaths, 2041--2050 & Models with more & Mean temperature, 2050 \\
    \midrule
    Austria     & $-9$\% & 7 of 24 & $-0.03$\,\textcelsius{} \\
    Germany     & $-3$\% & 9 of 24 & +0.13\,\textcelsius{} \\
    Czechia     & $-3$\% & 7 of 24 & +0.15\,\textcelsius{} \\
    Netherlands & $-10$\% & 8 of 24 & +0.18\,\textcelsius{} \\
    \bottomrule
  \end{tabular}

  \medskip
  \raggedright
  \begin{itemize}
    \item Fewer heat deaths, even though the year-round average temperature is
      about the same or slightly warmer.
    \item Possible reasons, none tested: summer heat differs in a way the
      yearly average hides; I only have temperatures for four sample years;
      or it's noise between models.
    \item Question: is it worth asking Caspar for summer temperatures?
  \end{itemize}
\end{frame}

\begin{frame}{Result 4: huge ratios where heat deaths are tiny}
  \small
  SSP5-8.5 relative to SSP2-4.5, heat deaths, 2041--2050:

  \medskip
  \centering
  \begin{tabular}{lcc}
    \toprule
    & Ratio & Heat deaths per year under SSP2-4.5 \\
    \midrule
    Greenland & 3.00 & 0.00001 \\
    Andorra   & 2.14 & 0.01 \\
    Monaco    & 2.14 & 0.11 \\
    \bottomrule
  \end{tabular}

  \medskip
  \raggedright
  \begin{itemize}
    \item These ratios are correct, but they mean little.
    \item \textbf{My suggestion:} keep every country, show deaths next to
      every ratio, and flag countries with very few deaths rather than drop
      them. Where to draw the line is a team call.
    \item Related, still open: how to mark countries with rare extreme
      events, such as Haiti (already marked) and Honduras.
  \end{itemize}
\end{frame}

\begin{frame}{Next steps}
  \begin{itemize}
    \item Done: the code for the combined summary, tested here.
    \item Caspar: one more run on his machine, for every year.
    \item Me, after his run: combined tables with uncertainty ranges, using
      whichever model choice the team prefers. Changing the model choice
      later will not need another run from Caspar.
  \end{itemize}
\end{frame}

\end{document}
